\part{Experiments, statistics and reproducibility}
\textit{Reconstructed section (the questionnaire was truncated; see the note in Part XII).}

\section{Experimental design}

\begin{qa}{E1}{Which experiments did you run?}
340 emulated runs in 8 studies (\code{configs/experiments/thesis.yaml}), plus 18 dataset-collection runs, real-OVS
runs and 5G-core runs:
\begin{center}\scriptsize
\begin{tabular}{@{}lll@{}}\toprule
Study & Design & Runs \\\midrule
main & 7 schedulers $\times$ 5 contexts $\times$ 5 seeds & 175 \\
ablation & 6 variants $\times$ 2 contexts $\times$ 3 seeds & 36 \\
capacity sensitivity & 2 schedulers $\times$ 2 contexts $\times$ \{50,100,150\} $\times$ 3 & 36 \\
load sensitivity & 2 $\times$ 2 $\times$ \{0.5,1,1.5\} $\times$ 3 & 36 \\
scalability & UEs per service \{1,2,4\} $\times$ 3 & 9 \\
resilience & 3 failures $\times$ 3 healing variants $\times$ 3 & 27 \\
timeline & 3 schedulers $\times$ 3 seeds (rural day) & 9 \\
in-loop SliceNet & 2 contexts $\times$ \{host, slicenet\} $\times$ 3 & 12 \\\bottomrule
\end{tabular}\end{center}
\end{qa}

\begin{qa}{E2}{Why 40\,s runs, 5\,s windows and a 1\,s tick?}
40\,s gives 8 windows per flow and 8 control intervals, enough for steady state after warm-up, and 340 runs fit in
about 1.5\,h with 4 workers. Resilience runs are 45\,s so that failure, recovery and repair all fit.
\honest{Long-run effects (minutes to hours: TCP fairness drift, diurnal patterns) are not captured.}
\end{qa}

\begin{qa}{E3}{Why 5 seeds in the main study and 3 elsewhere? Is that enough statistical power?}
The main study pairs 25 (context, seed) observations per comparison, which suits Wilcoxon. With the observed spread
(CIs of about $\pm$0.01--0.03 on utility), it can detect differences of about 0.02 or more. The observed allocator
differences ($<$0.012) are below that, so ``no significant difference'' means ``any difference is small'', not ``no
difference''.
\honest{3-seed studies are descriptive: CIs are reported, no tests.}
\end{qa}

\begin{qa}{E4}{How are seeds paired across schedulers?}
The same seed produces the same per-application traffic schedules for every scheduler in a context, so schedulers face
identical offered traffic, and comparisons are paired by (context, seed).
\end{qa}

\begin{qa}{E5}{Why the Wilcoxon signed-rank test rather than a t-test?}
Paired differences are not guaranteed to be normal (bounded metrics, heavy tails in delay). Wilcoxon is distribution-free
and robust, at a small loss of power.
\end{qa}

\begin{qa}{E6}{How do you handle multiple comparisons?}
Bonferroni over all scheduler $\times$ metric tests ($\alpha=\NBonferroniAlpha$). Both nominal and corrected significance are
reported. The healthcare-SLA effect of slicing survives correction ($p=\NSlaPNoSlicing$).
\end{qa}

\begin{qa}{E7}{Do you report effect sizes?}
Mean paired differences with CIs, in the metric's own units: utility, ms, shares.
\honest{No standardised effect sizes (Cohen's $d$ or the rank-biserial correlation); easy to add.}
\end{qa}

\begin{qa}{E8}{Were hypotheses fixed before seeing results?}
\code{docs/thesis/research_design.md} was written before the full matrix ran, and git history timestamps it. Verdicts are
computed mechanically by \code{experiments/report.py}.
\honest{The design and the evaluation were developed in the same project period, not externally pre-registered.}
\end{qa}

\begin{qa}{E9}{Which metrics are primary, and why?}
Primary: priority-weighted demand satisfaction (the allocation objective as stakeholders weigh it) and the
healthcare SLA-met share and p95 delay (critical-service protection). Secondary: fairness, loss, jitter, outage, MTTD
and MTTR.
\end{qa}

\begin{qa}{E10}{Isn't weighted utility circular, since the proposed scheduler is designed with those weights?}
Partly, yes. It is the stakeholders' objective, so a scheduler using it has an advantage by design.
That is why it is reported alongside unweighted metrics (Jain fairness, per-slice satisfaction, SLA-met share). Even
with this home advantage, the proposed scheduler is not significantly better, which makes the null result stronger,
not weaker.
\end{qa}

\begin{qa}{E11}{Why p95 delay rather than mean delay?}
Critical applications fail on tail delay (late ECG samples, frozen video), not on averages. p95 is computed per window
and summarised across windows; p99 is also reported.
\end{qa}

\section{Validity of the measurement campaign}

\begin{qa}{E12}{Were any runs invalidated? Why?}
Yes, twice, all archived and fully re-run:
\begin{enumerate}
\item \textbf{Busy-wait pacing} (111 runs): generators spun at about 100\% CPU each; host load reached 40--56 on 20
cores, perturbing kernel timers and therefore delays. Fixed with sleep-only pacing (\NGenCpuAfter\% CPU per 40\,Mb/s).
\item \textbf{Pre-OVS-validation} (340 runs):
  \begin{itemize}
  \item uplink-only utilisation (downlink saturation invisible);
  \item false failover under AQM-less queues;
  \item slow failover on OVS;
  \end{itemize}
  all discovered by running real OVS, fixed, and the matrix re-run.
\end{enumerate}
Archives: \code{results/runs_archive_*}.
\end{qa}

\begin{qa}{E13}{What other measurement bugs did you find?}
\begin{itemize}
\item \textbf{Partial final window}: traffic stops mid-window, and a stray near-empty last window divided delivered rate by
$N{+}1$ instead of $N$ (``static beats proposed at low load'' was this artefact). KPIs are now recomputed from each
run's dataset excluding the final window, uniformly for every run.
\item \textbf{Mislabelled RSNSD split}: the RSNSD baseline was reported as a ``group split'' when it was a random row
split (Q347).
\item \textbf{Hidden propagation delay}: netem on the bottleneck was silently replaced by HTB (delay placement redesigned).
\end{itemize}
\end{qa}

\begin{qa}{E14}{How do you know the measurement is right?}
\begin{itemize}
\item Delay floors equal the configured propagation to within 0.01\,ms (\NFloorEdge/\NFloorCore\,ms).
\item Unit tests feed synthetic captures with known one-way delay, loss and jitter and recover them exactly.
\item Probe RTT/2 agrees with the measured one-way delay.
\item All runs have packet-level data (\NRunsPcap/\NRuns).
\end{itemize}
\end{qa}

\begin{qa}{E15}{Could parallel runs interfere with each other?}
Each run has its own network namespace, capture directory and log directory, so packets and files cannot mix. Only
CPU is shared, and it is bounded by 4 workers (load stays below the 20 cores; no process spins).
\end{qa}

\begin{qa}{E16}{Are results from Linux bridges valid for Open vSwitch?}
The allocation and QoS mechanisms are identical (same tc qdiscs). Healing differs in timing: OVS failover took
\NOvsMttrAfter\,s and had a \NOvsOutageAfter\,s outage, against about 0.05\,s on bridges. Conclusions hold (sub-second
outage) but absolute numbers differ, and we report both.
\end{qa}

\begin{qa}{E17}{How reproducible is a single figure?}
Every figure is generated from \code{results/tables/*.csv} by \code{plot_thesis.py}; the tables come from per-run
artefacts via \code{aggregate.py}; each run is reproducible from its \code{run_meta.json} (parameters, seed, commit).
Bit-identical packet timing is not expected; distributions are.
\end{qa}

\begin{qa}{E18}{What does the full study cost to reproduce?}
About 1.5\,h for the matrix (4 workers, 20-core machine), minutes for training, under 1\,h for the 5G runs. No GPU.
About 100\,MB of captures per run (gitignored, regenerable).
\end{qa}

\section{Results by study}

\begin{qa}{E19}{Main study: what is the result?}
\begin{itemize}
\item Weighted utility \NUtilMin--\NUtilMax{} across the seven schedulers, with no significant difference.
\item Healthcare SLA-met share: \NHcSlaProposed{} (proposed) vs \NHcSlaNoSlicing{} (no slicing),
$\Delta=\NSlaDeltaNoSlicing$, $p=\NSlaPNoSlicing$.
\end{itemize}
Interpretation: \emph{enforcing} service slices with AQM matters; \emph{how} the shares are computed matters little
when borrowing is enabled.
\end{qa}

\begin{qa}{E20}{Ablation: what is the result?}
School:
\begin{itemize}
\item strict caps: utility \NAblSchoolFullUtil$\rightarrow$\NAblSchoolFullStrictUtil, healthcare p95
\NAblSchoolFullHc$\rightarrow$\NAblSchoolFullStrictHc\,ms;
\item no surplus redistribution (strict): $\rightarrow$\NAblSchoolNoSurplusStrictUtil;
\item no AQM: \NAblSchoolNoAqmUtil{}, \emph{higher}, because nothing is dropped; but general p95 \NAblSchoolFullGenP$\rightarrow$\NAblSchoolNoAqmGenP\,ms
and disaster healthcare SLA share \NAblDisasterFullHcSla$\rightarrow$\NAblDisasterNoAqmHcSla;
\item no context: \NAblSchoolNoContextUtil{} (identical to full);
\item oracle demand: \NAblSchoolOracleDemandUtil{} (an upper bound on what better demand estimation could add).
\end{itemize}
Borrowing is the dominant mechanism; surplus redistribution substitutes for it when borrowing is disabled; AQM trades
throughput for delay and SLA compliance; context awareness has no measurable effect.
\end{qa}

\begin{qa}{E21}{Sensitivity: what is the result?}
At 0.5$\times$ load all schedulers satisfy almost everything (static even edges ahead:
\NSensLoadSchoolStaticZeropFive{} vs \NSensLoadSchoolProposedZeropFive). At 1$\times$ the proposed scheduler leads
(\NSensLoadSchoolProposedOne{} vs \NSensLoadSchoolStaticOne); at 1.5$\times$ (\NSensLoadSchoolProposedOnepFive{} vs
\NSensLoadSchoolStaticOnepFive). H7 (``advantage everywhere'') is not supported. The advantage is conditional on
scarcity.
\end{qa}

\begin{qa}{E22}{Resilience: what is the result?}
\begin{itemize}
\item Fiber cut: outage \NOutCutFullMean\,s (full) vs \NOutCutNoneMean\,s (none).
\item Degradation: \NOutDegFullMean\,s (predictive) vs \NOutDegReactiveMean\,s (reactive) vs \NOutDegNoneMean\,s (none).
\item Core switch: \NOutCoreFullMean\,s in all variants; healing cannot fix it, but localises the fault in
\NCoreDiagnosedShare{} of runs.
\end{itemize}
\end{qa}

\begin{qa}{E23}{Timeline and scalability: what is the result?}
Timeline (a rural day with school and flood): all three schedulers adapt within one control interval of each context
change. Utility differences are within CIs (\code{timeline_summary.csv}). Scalability (1, 2, 4 UEs per service): the
pipeline and controller scale, and KPIs stay stable (\code{scalability_summary.csv}).
\honest{Scale is tiny compared with real deployments.}
\end{qa}

\begin{qa}{E24}{In-loop SliceNet: what is the result?}
Mixed-UE SLA-met share:
\begin{itemize}
\item school: \NMixedSchoolSlicenetSla{} (SliceNet) vs \NMixedSchoolHostSla{} (address-based); its education flows'
p95 fell from \NMixedSchoolHostEduPninefive{} to \NMixedSchoolSlicenetEduPninefive\,ms;
\item disaster: equal (\NMixedDisasterSlicenetSla{} vs \NMixedDisasterHostSla);
\item online classification accuracy \NMixedSchoolOnlineAcc.
\end{itemize}
H5b required gains in both contexts, so it is not supported. The benefit appears only where the mixed UE's misplaced
traffic actually contends.
\end{qa}

\begin{qa}{E25}{Which interactions between factors were not tested?}
\begin{itemize}
\item ablations combined (no AQM $\times$ strict caps);
\item failures during timelines or under overload;
\item RL policies under failures;
\item SliceNet under failures;
\item capacity $\times$ load jointly.
\end{itemize}
\end{qa}
